





























\title{On Graded Representations of Quantum Groups at Roots of Unity}
\author{Yang Jianwei}
\date{}
\begin{document}
\maketitle
\ifdefined\YTfallback
\noindent\fbox{\parbox{0.96\linewidth}{\small\textbf{[YangTex preview mode] }
The official class file \texttt{amsart.cls} was not found on this machine, so the
document uses a generic article layout. Put the class file into
\texttt{YangTex/templates/classes/} (or install it with \texttt{tlmgr}),
then recompile to obtain the real journal layout.}}
\bigskip
\fi
\begin{abstract}
We classify simple modules over a quantized enveloping algebra at a root of unity and describe the associated graded Grothendieck ring.
\end{abstract}
\paragraph{Keywords} quantum group; root of unity; graded representation; Grothendieck ring

\section{Introduction}
Let $U_q(\mathfrak{g})$ be the quantized enveloping algebra associated with a complex semisimple Lie algebra $\mathfrak{g}$ and let $q$ be a primitive $\ell$-th root of unity.

\begin{theorem}
The category of finite-dimensional graded $U_q(\mathfrak{g})$-modules is semisimple if and only if $\ell$ is coprime to the determinant of the Cartan matrix of $\mathfrak{g}$.
\end{theorem}

\begin{proof}
The proof reduces to the rank-one case by triangular decomposition and a careful analysis of the Steinberg-type tensor product formula.
\end{proof}

\section{Main estimates}
\begin{equation}
K_i E_j K_i^{-1}=q^{a_{ij}}E_j,\qquad [E_i,F_j]=\delta_{ij}\frac{K_i-K_i^{-1}}{q-q^{-1}},
\end{equation}
As a corollary we obtain a positivity result for the associated canonical basis.

\section{Concluding remarks}
The method extends to more general domains and boundary conditions; a detailed
treatment will appear elsewhere.

\begin{thebibliography}{9}
\bibitem{Evans} L.~C. Evans, \emph{Partial Differential Equations}, 2nd ed., AMS, 2010.
\bibitem{Brezis} H. Brezis, \emph{Functional Analysis, Sobolev Spaces and PDEs}, Springer, 2011.
\end{thebibliography}
\end{document}
